\section{Scope}

\subsection*{Release scope: v1 (Q3 2026)}

\begin{minipage}[t]{0.485\linewidth}
\begin{calloutbox}[height=8.6cm]{\faCheck\ In scope}
\footnotesize
\begin{itemize}
  \item Subscribe/unsubscribe to anomaly monitoring on any metric already in the Apex
    metrics catalog.
  \item Statistical anomaly detection (seasonal baseline + rolling standard deviation)
    on daily and hourly granularities.
  \item Delivery channels: in-app notification center, email, Slack (via existing Apex
    Slack integration).
  \item Per-metric sensitivity setting (low / medium / high).
  \item Maintenance-window muting.
  \item Mark-as-expected feedback loop that suppresses similar future alerts.
  \item Workspace-level subscription audit view for admins.
\end{itemize}
\end{calloutbox}
\end{minipage}\hfill
\begin{minipage}[t]{0.485\linewidth}
\begin{riskbox}[height=8.6cm]{\faTimes\ Out of scope}
\footnotesize
\begin{itemize}
  \item Custom SQL / ad-hoc metric alerting outside the standard catalog.
  \item Forecast-based (``will breach in N days'') predictive alerting.
  \item PagerDuty, Opsgenie, or other third-party incident-tool routing.
  \item Admin-configured default subscriptions on behalf of other users.
  \item Mobile push notifications (mobile app does not yet exist).
  \item Alert-triggered automated actions (e.g.\ auto-rollback, auto-ticket creation).
\end{itemize}
\end{riskbox}
\end{minipage}
\par

\subsection*{Explicitly deferred to v2}
Sub-minute detection latency, cross-metric correlation alerts (``these three metrics
moved together''), and a public alerting API for customers to build their own
integrations are all validated needs but are sequenced after v1 usage data confirms the
core loop works.
